\section{Introduction}
\label{sec:introduction}

Blockchain technology has emerged as a promising foundation for decentralized and auditable healthcare data sharing~\cite{kakarlapudi2021systematic, welzel2025enabling, li2024you}.
Its transparency and immutability allow consent enforcement and authorization decisions to be verified and traced across institutions.
However, the same transparency means that consent enforcement on a public blockchain can expose information about the governing consent policies, even when the health data themselves remain encrypted~\cite{beres2021blockchain, li2025privacy}.

This exposure matters because consent policies in healthcare are often privacy-sensitive~\cite{lawson2021data}.
A \emph{consent policy} encodes the conditions under which a data subject authorizes others to access or use their health data, such as who may access the data, for what purpose, and under what constraints.
For example, a policy based on biomedical data-use terms may reveal whether a dataset is restricted to disease-specific research, whether ethics approval is required, or whether access is limited to particular institutions or projects.
If such a policy is written into a smart contract, embedded in verification logic, or reflected in public execution traces, unauthorized observers may infer information about the protected resource, the data subject's constraints, or the requester's authorization context.

Existing privacy-preserving blockchain solutions do not address this problem, because they mainly protect data contents, user identities, or transaction details~\cite{rondelet2019zeth, williamson2018aztec, nadler2023tornado}.
Encrypted data-sharing mechanisms such as attribute-based encryption (ABE) and proxy re-encryption (PRE) control access to data~\cite{ibraimi2009ciphertext, ullah2023survey}, and zero-knowledge transaction systems hide asset ownership, values, or participant relationships~\cite{zhou2024leveraging}; none of them protects the confidentiality of the policies that govern access decisions.
Hidden-policy ABE~\cite{yang2024hamed,ding2026badsabe} does hide policy content, but it does not control what a public chain exposes while it enforces the policy.
On a public chain, a policy can leak not only through explicit disclosure but also through metadata: the size of the calldata, the layout of the public inputs, the emitted events, the verifier interface, and the gas consumed.
Policy-private consent enforcement therefore requires more than hiding data or proving correctness in zero knowledge: it must also control this metadata leakage, that is, what the public-chain transcript (the transaction data, events, and gas visible to any observer) reveals about the hidden policy.

This paper studies this problem for blockchain-based healthcare data sharing on public blockchains.
The challenge is to let an on-chain verifier check that a requester satisfies a provider-defined consent policy, while ensuring that observers who see only public-chain metadata learn minimal information about the governing policy and the requester's authorization evidence.
We focus on Boolean consent policies, including policies derived from biomedical data-use conditions such as terms from the Data Use Ontology (DUO), the Global Alliance for Genomics and Health (GA4GH) standard vocabulary of consent-based data-use conditions~\cite{lawson2021data}, and disease-scope constraints~\cite{WHO_ICD10}.

To address this problem, we propose the Trusted Proof Service-assisted Zero-Knowledge Policy Interpreter (TPS-ZKPI).
In TPS-ZKPI, the data provider commits to a hidden consent policy on-chain and provisions the corresponding policy-bound proving material to a Trusted Proof Service (TPS).
A requester submits its policy-relevant authorization evidence, such as attributes, credentials, and access context, to the TPS, which constructs the private witness and generates a zero-knowledge proof that the requester satisfies the committed policy.
A smart contract then verifies this proof against the provider's policy commitment.
As a result, every authorization decision remains publicly auditable, while the policy logic never appears in the public-chain transcript.

To prevent the transcript itself from revealing the policy, TPS-ZKPI combines four design elements.
First, proof generation at the TPS keeps confidentiality-sensitive policy material and requester evidence off the chain.
Second, heterogeneous Boolean consent policies are compiled into fixed-length Reverse Polish Notation (RPN) programs, so that every policy has the same length.
Third, building on the standard commit-and-prove paradigm, a reusable zero-knowledge policy interpreter evaluates these programs; this avoids policy-specific verifier contracts and allows policies to evolve without changing the public verification interface.
Fourth, transcript normalization, including randomized gas padding, reduces the remaining policy-dependent variation in calldata, event logs, public parameters, and gas consumption.

We evaluate TPS-ZKPI against a passive public-chain observer who sees only the public-chain transcript.
The evaluation expresses consent policy privacy as three objectives: Policy Content Hiding (PCH), Policy Indistinguishability (PI), and Structure and Parameter Privacy (SPP).
For these objectives, supervised and unsupervised attacks test whether the public-chain metadata allow the observer to recover hidden policy content, to distinguish otherwise-equivalent policies, or to infer structural and parameter properties of a policy.
The results show that $2$ of the $45$ supervised tasks meet our leakage criterion, which flags a task only when its mean AUC exceeds $0.55$ and the attack is statistically significant in at least half of $30$ policy-disjoint splits; both failures trace to chance correlations with the randomly drawn padding targets.
Furthermore, a gas-padding ablation shows that, without padding, the same attacks recover a policy-correlated gas signal, which randomized gas padding suppresses.

This paper makes the following contributions:

\begin{itemize}
	\item We propose TPS-ZKPI, a TPS-assisted zero-knowledge approach to publicly
	      verifiable and auditable consent enforcement. The soundness of on-chain
	      verification and the auditability of submitted decisions hold without
	      trusting the TPS, whereas policy confidentiality, faithful witness
	      construction, and liveness rely on trusting the TPS, which keeps policy-bound proving
	      material and requester authorization evidence off-chain.

	\item We encode each Boolean consent policy as a fixed-length RPN program over a bounded token alphabet and evaluate it in a single zero-knowledge circuit.
	      Combined with policy commitments, it gives all policies the same
	      policy-independent on-chain verification interface: one circuit serves
	      every policy, and policies can evolve without changing the on-chain verifier.

	\item We empirically evaluate how much policy information remains in the
	      normalized public-chain metadata on the Polygon~Amoy
	      testnet\footnote{\url{https://amoy.polygonscan.com/txs?a=0x461853d10fB9A286a75dE05c9B166e77645160D7}}, using supervised
	      classification attacks judged by our leakage criterion and unsupervised
	      clustering attacks. A paired gas-padding ablation further shows
	      that randomized gas padding is responsible for suppressing the modeled
	      gas channel and quantifies its gas overhead.
\end{itemize}

The remainder of this paper is organized as follows.
Section~\ref{sec:background} reviews the technical foundations: cryptographic access control in blockchain healthcare, zero-knowledge proofs, and the Data Use Ontology.
Section~\ref{sec:related-work} discusses related work and positions TPS-ZKPI against existing blockchain consent-management, privacy-preserving blockchain, hidden-policy, private policy-evaluation, reusable proving, and metadata-leakage approaches.
Section~\ref{sec:problem} defines the system and trust model, the public-chain adversary, and the privacy objectives.
Section~\ref{sec:design} presents the TPS-ZKPI design.
Section~\ref{sec:eval} reports the empirical evaluation.
Section~\ref{sec:discussion} discusses limitations and future work, and
Section~\ref{sec:conclusion} concludes the paper.
